% ============================================================
% 符号表（E-M2）· 可直接 \input 进论文的 LaTeX 片段
% ------------------------------------------------------------
% 用途 / Purpose ：全队统一符号的单一出口——建模、代码、论文用同一套符号。
% 对应工作流 / Workflow ：workflows/03-问题分析与假设.md §符号说明 ·
%                          workflows/04-建模.md §第三步（定义变量、参数）·
%                          workflows/05-编程求解.md §代码规范
% 门禁指针 / Gates ：Euler-ENGINE.md §1 E-M2——符号与代码/论文一致；§3「建模」行
% ------------------------------------------------------------
% 用法 / Usage ：
%   1) 论文「符号说明」节里写 \section{符号说明} 后接 \input{sections/符号表.tex}
%      （节标题由论文给，本片段不含 \section；如需自带标题，启用下方注释行后删掉论文里的 \section）
%   2) 包依赖：longtable + booktabs + tabularx + array
%      —— 模板 cls：美赛-MCM-ICM/mcmthesis.cls 随包；国赛-CUMCM/cumcmthesis.cls 需自备（见该目录 README）；
%      换用其他模板若缺包，在导言区加 \usepackage{longtable,booktabs,tabularx,array}
%   3) 占位符 <…> 仅在「含义/首次出现处/单位」等**文本模式**列内使用；交稿前必须全部替换
%      （pdflatex+T1 下 < > 会渲染成 ¡ ¿，属占位现象，不影响逻辑）
%   4) **数学模式内不要写中文**（会掉字/tofu）：符号列占位故意用 ASCII 记号 $\mathrm{sym}_N$ 演示；
%      真实符号请替换为 $x_i$、$\beta$、$\lambda$ 等。中文一律写进「含义」列（文本模式）。
%      确需在数学环境里夹中文，用 \text{中文} 包（依赖 amsmath，双模板已载）。
%   5) 表体用 longtable（跨页自动续排）；符号较少（≤15 行）想固定位置，改用文末 tabularx 变体
%   6) 单位口径：有单位写 SI 单位（如 $m/s$、$kg\cdot m^{-2}$）；无量纲写 $1$ 或 --
% ============================================================

% \section{符号说明}   % ← 需要自带节标题时启用本行

\begin{longtable}{c l c c}          % 符号 | 含义 | 单位 | 首次出现处
  \caption{符号说明}\label{tab:symbols}\\
  \toprule
  \textbf{符号} & \textbf{含义} & \textbf{单位} & \textbf{首次出现处} \\
  \midrule
  \endfirsthead
  \multicolumn{4}{r}{\small 续表~\thetable}\\
  \toprule
  \textbf{符号} & \textbf{含义} & \textbf{单位} & \textbf{首次出现处} \\
  \midrule
  \endhead
  \midrule
  \multicolumn{4}{r}{\small 续下页}\\
  \endfoot
  \bottomrule
  \endlastfoot
  % ---- 主符号（按首次出现顺序排列；符号列替换为真实符号，勿在数学模式写中文）----
  $\mathrm{sym}_1$ & <含义：中文全角说明，首次出现附英文术语> & <单位 1>            & <式(1)> \\
  $\mathrm{sym}_2$ & <…>                                     & <单位 2>            & <式(3)> \\
  $\mathrm{sym}_3$ & <…>                                     & <无量纲（写 $1$）>    & <第 5.1 节> \\
  $\mathrm{sym}_4$ & <…>                                     & <单位 4>            & <步骤 2> \\
  % ---- 下标 / 集合 / 算子 ----
  $\mathrm{sym}_5$ & <…>                                     & <单位 5 / 无量纲>    & <…> \\
  % ---- 参数（带标定来源）----
  $\mathrm{sym}_6$ & <…（取值来源：实测/文献/拟合）>            & <单位 6>            & <表 2> \\
\end{longtable}

% ------------------------------------------------------------
% 【变体】符号少且需固定位置的 tabularx 版（与上面二选一，不要同时启用）
% \begin{table}[htbp]
%   \centering
%   \caption{符号说明}\label{tab:symbols}
%   \begin{tabularx}{\textwidth}{c X c c}
%     \toprule
%     \textbf{符号} & \textbf{含义} & \textbf{单位} & \textbf{首次出现处} \\
%     \midrule
%     $\mathrm{sym}_1$ & <含义> & <单位 1> & <式(1)> \\
%     \bottomrule
%   \end{tabularx}
% \end{table}
% ------------------------------------------------------------

% ============================================================
% 【一致性检查卡】符号 ↔ 代码变量 ↔ 论文 —— 三处一致清单
% 本卡为 LaTeX 注释，不参与编译；在**本工作副本**上逐条打勾（复制到赛期工作副本再填）。
% 检查时机：① 每轮模型迭代后 ② 代码改动后 ③ 论文定稿前（E-M7 四轮自审第三轮）
% ------------------------------------------------------------
% 表：逐符号登记
% | 符号 | 代码变量名（脚本:行） | 论文出现处（式/表/图） | 单位一致 | 三处一致 |
% |------|----------------------|----------------------|---------|---------|
% | <符号1> | <beta（solve_q1.py:42）> | <式(3)/图2/表1> | [ ] | [ ] |
% | <符号2> | <…> | <…> | [ ] | [ ] |
% | <符号3> | <…> | <…> | [ ] | [ ] |
%
% 判据（逐条打勾，任一条不过 → 就地修正并重核）：
% [ ] 符号表内每个符号**有且仅有一行**（无重复、无别名）
% [ ] 论文正文用到的每个符号**都在符号表里**（无未登记符号）
% [ ] 符号表里的每个符号**都在正文被使用过**（无僵尸符号；仅附录用到的另加注）
% [ ] 代码变量名与符号**一一对应**（同一物理量不得出现两个变量名；缩写/别名全清）
% [ ] 单位四处一致：符号表 · 坐标轴标签 · 表头 · 正文（Euler workflows/07 §格式细节）
% [ ] 量纲检查通过：方程/目标函数两边量纲一致（Euler workflows/05 §结果验证）
% [ ] 写法统一：$\cdot_i$ 与 $\cdot_{i}$、希腊字母大小写、加粗向量写法全文一致
% [ ] 单位制统一（SI 或题给单位），换算系数已在符号表或正文注明
% [ ] 论文图注措辞与符号表一致（workflows/07 §图表规范）
% [ ] 参数取值来源已登记（实测/文献/拟合/量级论证），与假设表「依据」列口径一致
%
% 不过怎么办：
%   1) 不要手改 PDF/最终稿；回到源头（代码变量名 / 论文源文件 / 本表）改一处；
%   2) 改动波及 → 同步「假设表 · 对应环节」「模型卡 · 假设链接」「受影响图表」三处；
%   3) 论文定稿后新增符号 → 重新编译并重跑本卡（禁只在 PDF 上加注）。
% ============================================================
